\section{The mathematical definition}\label{sec:09-rings:01-definition:1}

A \textbf{ring} is a set \(R\) with \emph{two} binary operations, addition and multiplication, such that:

\[
\begin{aligned}
\text{(R1)}&\quad (R, +, 0, -(-)) \text{ is a commutative group} \\
\text{(R2)}&\quad (a \cdot b) \cdot c = a \cdot (b \cdot c) \quad \text{(multiplication is associative)} \\
\text{(R3)}&\quad \exists\, 1 \in R,\ 1 \cdot a = a \text{ and } a \cdot 1 = a \quad \text{(multiplicative identity)} \\
\text{(R4)}&\quad a \cdot (b + c) = a \cdot b + a \cdot c, \quad (a + b) \cdot c = a \cdot c + b \cdot c \quad \text{(distributivity)}
\end{aligned}
\]

Some textbooks do not require a multiplicative identity. Such a structure is called a \emph{rng}. The missing ``i'' stands for the missing identity. This book includes the identity, since that is the more common convention.

Note that (R1) requires \textbf{commutative} addition, unlike the general \texttt{Group} of Chapter 7. Hence, before defining \texttt{Ring}, we first define what ``commutative'' means as an extension of \texttt{Group}.

\subsection{Sources, quoted}\label{sec:09-rings:01-definition:2}

Formal definitions and citations for this section, gathered here for reference (full entries in the \hyperref[sec:root:bibliography:1]{Bibliography}):

\begin{itemize}
\tightlist
\item
  \textbf{Ring.} A set with two binary operations, addition and multiplication, such that the set is an abelian group under addition, multiplication is associative and distributes over addition, and (in the convention this book follows) a multiplicative identity exists (\cite{DummitFoote2003}, Ch. 7 ``Introduction to Rings,'' §7.1 ``Basic Definitions and Examples''). Dummit and Foote's own base definition does \emph{not} require a multiplicative identity, calling that weaker structure a ring ``not necessarily with 1'' and reserving ``ring'' by default for the identity-including case in most of the book's later chapters, the same choice (R3) above makes explicit. This is a structural citation to the section and its numbered axioms, not a verified word-for-word excerpt.
\end{itemize}
